\documentclass[11pt]{article}
% Shared preamble for the bite-sized LaTeX doc series (docs/latex/*.tex).
% \input this at the top of each numbered file rather than duplicating the
% package list -- keeps every file's own content close to (and comfortably
% under) 10 pages.
\usepackage[a4paper,margin=2.5cm]{geometry}
\usepackage{amsmath,amssymb}
\usepackage{hyperref}
\usepackage[backend=biber,style=numeric,sorting=none]{biblatex}
\usepackage{listings}
\usepackage{xcolor}
\usepackage{enumitem}
\addbibresource{references.bib}

\lstdefinelanguage{Rust}{
  keywords={fn,let,mut,pub,struct,enum,impl,match,if,else,for,while,return,use,mod,self,Self,const,static},
  keywordstyle=\color{blue}\bfseries,
  comment=[l]{//},
  commentstyle=\color{gray}\itshape,
  string=[b]",
  stringstyle=\color{teal},
  basicstyle=\ttfamily\small,
  showstringspaces=false,
  breaklines=true,
}
\lstset{language=Rust}

\hypersetup{colorlinks=true, linkcolor=blue, citecolor=blue, urlcolor=blue}


\title{The Non-Analytic Critical-Region Term \\ \large Part 3: Rust Implementation}
\author{outram-park-fork-coolprop}
\date{2026-07-10}

\begin{document}
\maketitle

\begin{abstract}
Parts 1--2 covered the theory and the numerical guard. This document covers
how the non-analytic term is represented as data, dispatched, and evaluated
in this crate's Rust code — the enum-dispatch pattern this workspace uses
everywhere in place of trait objects, and how a fluid's hardcoded data feeds
into it.
\end{abstract}

\section{Where the data lives: an enum variant, not a class}

Following this workspace's mandatory ``no trait objects'' rule, every
residual-Helmholtz term \emph{type} CoolProp supports is a variant of a
single enum, \texttt{ResidualTerm} (\texttt{src/eos.rs}):
\begin{lstlisting}[language=Rust]
pub enum ResidualTerm {
    Power { n: &'static [f64], d: &'static [f64], t: &'static [f64], l: &'static [f64] },
    Gaussian { n: &'static [f64], d: &'static [f64], t: &'static [f64],
               eta: &'static [f64], epsilon: &'static [f64],
               beta: &'static [f64], gamma: &'static [f64] },
    NonAnalytic { n: &'static [f64], a: &'static [f64], b: &'static [f64],
                  beta: &'static [f64], big_a: &'static [f64], big_b: &'static [f64],
                  big_c: &'static [f64], big_d: &'static [f64] },
    // ... Exponential, DoubleExponential, GaoB
}
\end{lstlisting}
Each field is a \texttt{\&'static [f64]} slice, not a \texttt{Vec} — the
whole term list for a fluid is \texttt{const} data compiled directly into
the binary (see Part 4's \texttt{water.rs}), never read from a file at
runtime. A fluid's full equation of state, \texttt{FluidEos}, is simply
\texttt{residual: \&'static [ResidualTerm]} plus the ideal-gas terms and
reducing/critical parameters.

\section{Dispatch: an exhaustive \texttt{match}, not virtual calls}

Evaluating a fluid's residual Helmholtz derivatives at a state
$(\delta,\tau)$ walks the term list and dispatches on the enum:
\begin{lstlisting}[language=Rust]
pub fn accumulate(&self, delta: f64, tau: f64, acc: &mut HelmholtzDerivs) {
    match self {
        ResidualTerm::Power { n, d, t, l } => { /* ... */ }
        ResidualTerm::Gaussian { n, d, t, eta, epsilon, beta, gamma } => { /* ... */ }
        ResidualTerm::NonAnalytic { n, a, b, beta, big_a, big_b, big_c, big_d } => {
            accumulate_non_analytic(n, a, b, beta, big_a, big_b, big_c, big_d, delta, tau, acc);
        }
        // ...
    }
}
\end{lstlisting}
This is the enum-dispatch pattern the whole workspace uses in place of
\texttt{Box<dyn Trait>}: the compiler checks the \texttt{match} is
exhaustive, so adding a new CoolProp term type without updating this
function is a compile error rather than a silent runtime gap — and there is
no heap allocation or virtual-call indirection, since every branch is known
at compile time.

\section{The non-analytic evaluator itself}

\texttt{accumulate\_non\_analytic} takes the eight coefficient slices, the
state $(\delta,\tau)$, and a mutable accumulator, applies the branch-point
offset guard from Part 2, then loops over however many non-analytic terms
the fluid has (Water has two) accumulating each term's contribution:
\begin{lstlisting}[language=Rust]
for i in 0..n.len() {
    let (ni, ai, bi, betai) = (n[i], a[i], b[i], beta[i]);
    let (aai, bbi, cci, ddi) = (big_a[i], big_b[i], big_c[i], big_d[i]);
    let dm1 = delta - 1.0;
    let dm1_sq = dm1 * dm1;

    let theta = (1.0 - tau) + aai * dm1_sq.powf(1.0 / (2.0 * betai));
    // ... dtheta_ddelta, d2theta_ddelta2, psi and its derivatives,
    // big_delta and its derivatives, delta_bi = big_delta.powf(bi)
    // and ITS derivatives -- see Part 2's derivative chain --

    acc.add_scaled(&HelmholtzDerivs {
        a: delta * ni * delta_bi * psi,
        ad: /* ... */, at: /* ... */,
        add: /* ... */, att: /* ... */, adt: /* ... */,
    });
}
\end{lstlisting}
\texttt{acc.add\_scaled} sums this term's contribution into the running
total for the whole residual sum — every term type ultimately contributes
through the same six-field \texttt{HelmholtzDerivs} accumulator, which is
what lets \texttt{state\_trho} (Part 4) stay agnostic to which term types a
particular fluid happens to use.

\section{Why this design, not something more ``generic''}

A tempting alternative would be a trait \texttt{ResidualTermKind} with one
implementor struct per term type, dispatched via \texttt{Box<dyn
ResidualTermKind>}. This workspace deliberately avoids that (see the
root \texttt{CLAUDE.md}'s ``no trait objects'' rule) for three concrete
reasons that matter for a physics library specifically: (1) the set of term
types is closed — CoolProp's own generator (Part 4) already aborts loudly
if a fluid ever needs an unimplemented type, so exhaustiveness-checking via
\texttt{match} is strictly more valuable than open-ended extensibility;
(2) every fluid's entire EOS is \texttt{const} data with zero heap
allocation, which a \texttt{Vec<Box<dyn Trait>>} would break; (3) go-to-definition
on an enum variant works reliably in any IDE, while go-to-definition on a
\texttt{dyn Trait} call site often does not — directly serving this
workspace's ``human interface layer'' design principle (also in the root
\texttt{CLAUDE.md}) that any public API must be navigable by
\texttt{rust-analyzer} alone.

\printbibliography
\end{document}
